% !TeX root = ../papers/nc-p3.tex
\section{二変数のSmith指数と外側の極の全域}
\label{sec:rank-fifteen-outer-pole-fan}
\begin{lemma}\label{lem:rank-fifteen-two-parameter-smith}
$R$ を $\mathbb Z_{(3)}[S,T]_{(3,S,T)}$ のHensel化とし、$M$ を $R$ 上の行列とする。
高さ一の素点 $(S),(T)$ で必要な本数までのSmith指数をそれぞれ昇順に
$a_j,b_j$ とする。剰余の対角線 $S=T=t$ 上の $\mathbb F_3[[t]]$ で
同じ本数の指数が $c_j=a_j+b_j$ なら、対角線の主元を用いて
$R$ 上で単項式 $S^{a_j}T^{b_j}$ を対角成分として順に分離できる。
\end{lemma}
\begin{proof}
$R$ は正則局所環で、$S,T$ は非同伴な素元である。全成分の両素点での付値が
$a_1,b_1$ 以上なので、$S^{a_1}T^{b_1}$ は全成分を割る。
この単項式で割った行列は対角線で単元主元を持ち、その主元は $R$ でも単元である。
これを消去すると、両素点での次の指数は $a_j-a_1,b_j-b_1$、
対角線では $c_j-c_1$ であり、同じ等式が続く。帰納し、分離の際に除いた
共通単項式を掛け戻せばよい。特に対角線で選んだ行・列と同じ語基底が使える。
\end{proof}

\begin{theorem}\label{thm:rank-fifteen-outer-pole-fan}
通常な標数零の階数15窓で、剰余標数3の付値について
$0<2\nu(x-1)<\nu(y)$ とする。この外側の領域は全て標準Koszul左右AS正則で、
各始点で四中心四次式上自由階数256である。
\end{theorem}
\begin{proof}
補題\ref{lem:rank-fifteen-outer-pole}の $(w,u)$ はともに正付値である。
$(\nu(w),\nu(u))$ の第一象限を、順に
\[
          (1,0),\quad(1,1),\quad(1,2),\quad(0,1)
\]
の隣り合う二辺で分ける。最初と最後は付値零の境界で、内部の二つの半直線だけを
別に扱う。$(1,1)$ は定理\ref{thm:rank-fifteen-outer-pole}による。

$(1,2)$ では $u=mw^2$, $s=w^2$ として $H^2=m^2s^3$ である。
$B=(x-1)/s$ の剰余は $-4$、$y/(x-1)^6$ の剰余は $m^4/4096$ であり、
前節と同じ単元ヤコビ行列によって所定の $x,y$ に持ち上がる。
$\mathbb Q(m)$ 上で四十八桁を使うと、未切断経路は $(0,4,4,6,6)$、
四次中心は $(0,3,3,3)$、五次の積は五桁を要する。長さ6から13では
$(5,2,3,2,4,0,3,0)$ 桁を使う。実際に用いた全ての主元の定数分子・分母の
剰余の平方自由な積は $m$ であり、従ってこの半直線の全ての単元に適用できる。

隣り合う二辺を $(a,b),(c,d)$ と書くと $ad-bc=1$ なので、開いた錐の内部では
$w=S^a T^c$, $u=S^b T^d$ と一意に書け、$S=w^du^{-c},\ T=w^{-b}u^a$ はともに正付値である。
全平方根は定数項を指定した単純根なので、経路係数は補題の $R$ に入る。
$(1,0)$ に沿う一般指数は $w$ を消す接方向のものである。実際
$HL=w$, $H=uw$, $z=(1+\sqrt{1-u^2})/u$ によって $z$ は境界の一般元となる。
$(0,1)$ に沿う一般指数は $u$ を消す非分岐極のもので、
$A=(\sqrt{1-w^2}-w)/\sqrt2$, $B=(\sqrt{1-w^2}+w)/\sqrt2$ は非定数の円錐曲線の表示である。
それぞれ一様化元は既存表示の $H$ または $r$ の単元倍であり、
第\ref{sec:rank-fifteen-unit-pole-depth}節の両指数列が使える。
内部の二辺ではそれぞれ $u=mw,\,u=mw^2$ における一般の $m$ を使う。
その一般指数は $m=1$ の剰余でのものと一致する。実際それを基準に選んだ各消去は、
単元主元が尽きた際の定数項の消滅を $\mathbb Q(m)$ 上で厳密に確認している。

\path{computations/verify_rank_fifteen_outer_pole_fan.py} は各錐の対角線、すなわち
$(w,u)=(t^2,t),(t^2,t^3),(t,t^3)$ でそれぞれ48、79、64桁の標数零計算を行う。
各錐の未切断経路80組、四次中心4組、五次までの中心商96組、長さ6から13の左右の積の照合128組で、
対角線の指数は両辺の昇順の指数の項別の和に一致する。比較は各行列の全主元に対して行う。

未切断経路から順に補題を適用し、対角線の語基底で自由格子と積の表示を作る。
各辺の離散付値環上では、既知の基底とこの基底が同じ経路格子を張るので、
基底変更は可逆である。中心の自由座標を指定する最後の始点の全四次語基底と列番号も
両辺と対角線で一致する。従って同じ四本の中心を課して帰納できる。
左右の積の照合では、中間の語基底の変更も行と列の可逆変換に過ぎない。
十三次まで各段の一般の次元下界に達する本数を消去でき、
十三次で商は零である。所定の窓で $S,T$ はともに非零なので各単項式での除算が使え、
従来の自由加群と左右の解消の判定を適用できる。
\end{proof}
